\section{Discussion}

\subsection{Synthesis of Findings}
The results demonstrate that eHealth systems adoption is not determined by isolated technical capabilities, but by the alignment of user readiness, technical interoperability, institutional frameworks, and organizational execution. Across the systematic review and case analyses, these domains consistently interact to shape adoption outcomes.

\subsection{Cross-Case Comparative Insights}
The comparative analysis highlights important differences in how adoption determinants manifest across implementation contexts. Estonia represents a highly aligned system, where institutional mandate, technical infrastructure, and organizational execution converge. In contrast, the VA/DoD case illustrates how misalignment—particularly in user satisfaction and implementation governance—can undermine adoption despite strong institutional backing.

These findings suggest that institutional mandate alone is insufficient. Effective adoption requires coordinated execution across all domains, particularly at the organizational and user levels.

\subsection{Implications for Interoperability Strategy}
The findings reinforce that interoperability is not solely a technical problem. While standards and interfaces are necessary, they must be supported by governance frameworks, workflow integration, and organizational alignment.

Policies such as the HITECH Act and the 21st Century Cures Act attempt to address these challenges by realigning incentives and reducing information blocking. However, implementation effectiveness depends on how these policies translate into organizational and technical practice.

\subsection{Theoretical Perspective: Market Dynamics and Ethical Alignment}

The adoption of eHealth systems occurs within a complex socio-technical and economic environment. Healthcare IT markets are characterized by vendor competition and proprietary system architectures that incentivize innovation but also produce structural fragmentation across EHR systems and HIEs. From a systems perspective, individual firms optimizing for competitive advantage—through proprietary data models and ecosystem lock-in—generate outcomes that are suboptimal at the system level, where continuity of information is essential for safe care delivery. Institutional mechanisms such as the HITECH Act, the 21st Century Cures Act, and information blocking regulations represent deliberate attempts to correct this misalignment by realigning incentives toward system-wide interoperability and data sharing.

Beyond policy and economics, interoperability may be understood not only as a technical requirement but as an ethical imperative: it supports continuity of care, reduces medical error, and promotes equitable access to health information. Sustainable eHealth adoption requires harmonization between technological capability, market incentives, regulatory frameworks, and ethical principles—an alignment that the U-T-I-O Model is designed to facilitate.


\subsection{Research Questions Revisited}

The findings directly address each of the three research questions guiding this study. \textit{RQ1 (What determinants most consistently influence eHealth systems adoption?)} is answered by the four-domain synthesis of the U-T-I-O framework: user readiness, technical interoperability, institutional alignment, and organizational execution collectively and consistently predict adoption outcomes across both the literature and the case analyses. \textit{RQ2 (How are these determinants reflected in real-world implementations?)} is addressed through the comparative case analysis, which demonstrates that the relative strength of each domain is traceable to specific design and governance decisions in each implementation context—Estonia's legal mandate and X-Road architecture, the NHS's centralized interoperability vision, the ONC's market-based incentive model, and the VA/DoD's procurement and execution challenges. \textit{RQ3 (How can these determinants be synthesized into a reusable adoption model?)} is answered by the U-T-I-O Model itself, which provides a framework that is grounded in empirical evidence, theoretically integrated with established IS and innovation adoption literature, and operationalizable across diverse healthcare contexts.

\subsection{Practical Implications}
For healthcare leaders and system architects, the results suggest that adoption strategies should prioritize:

\begin{itemize}
\item Alignment between technical architecture and clinical workflows, ensuring that interoperability investments translate directly into usability improvements at the point of care
\item Strong governance and leadership engagement throughout implementation, not only at launch but continuously across the adoption lifecycle
\item Investment in user training, satisfaction monitoring, and feedback mechanisms, treating user readiness as an ongoing operational concern rather than a one-time onboarding activity
\item Compliance with interoperability standards and regulatory frameworks as a baseline, while actively leveraging policy incentives (HITECH, Cures Act, information blocking rules) to align vendor and organizational behavior toward system-wide interoperability goals
\end{itemize}

Organizations that address only technical interoperability without equivalent investment in organizational and user readiness domains are unlikely to achieve sustained adoption, regardless of the strength of their institutional mandate or the quality of their technical infrastructure.